\documentclass[10pt,a4paper]{amsart}
\usepackage[margin=19mm]{geometry}
\usepackage{amsmath,amssymb,mathtools}
\usepackage[scr=rsfs,cal=euler]{mathalfa}
\usepackage{xfrac}
\usepackage[unicode,hidelinks,bookmarks=false]{hyperref}
\newcommand{\ie}{i.e.\ }
\newcommand{\cf}{cf.\ }
\newcommand{\resp}{respectively\ }
\newcommand{\Subsection}{\subsection}
\providecommand{\nobreakdash}{\nobreak\hbox{-}}
\setlength{\emergencystretch}{2em}
\multlinegap0pt
\usepackage{amssymb}
\usepackage{mathtools}
\usepackage[shortlabels]{enumitem}
\newtheorem{lemma}{Lemma}
\newcommand{\nn}{\nonumber}
\newcommand{\norm}[1]{\left\lVert #1\right\rVert}
\DeclareMathOperator{\sinc}{sinc}
\title{Counterexamples to Whole-Sequence Convergence of Variable-Smoothing Full-Splitting Methods}
\author{Min Tao}
\date{Source: arXiv:2608.20859v3 (CC BY 4.0)}
\begin{document}
\maketitle

\noindent\emph{Source and licence.} Counterexamples to Whole-Sequence Convergence of Variable-Smoothing Full-Splitting Methods. Version arXiv:2608.20859v3 (CC BY 4.0), distributed by arXiv under the Creative Commons Attribution 4.0 International licence, http://creativecommons.org/licenses/by/4.0/, as recorded in the arXiv licence metadata for this exact version and shown on the abstract page https://arxiv.org/abs/2608.20859v3. Full licence notice: \texttt{licenses/arxiv-2608.20859-CC-BY-4.0.txt}.\par\smallskip

\noindent\textit{Editorial scope.} Counted unit: the article's lemma lem:elementary (source0.tex lines 2913--2933). The article's complete original proof of it is delivered with it (source0.tex lines 2934--3375, one \texttt{proof} environment of 442 lines). No mathematics is written, repaired or re-derived: the statement and the proof are the article's own lines, in the article's own notation, and no retained line is substituted or re-flowed.  Retained uncounted as the source-local context the proof needs: lines 1462--1467; lines 2853--2860.  Nothing was omitted from any retained span, so the package carries no omission record.  No retained line contains a citation, so the wrapper carries no bibliography and no bibliography entry.  No retained line contains a figure environment, an image-inclusion command or drawing code, so no image is delivered and the package declares no asset dependency.  Editorial formatting only, in addition: an AMS \texttt{amsart} preamble that restates the article's own declarations and macros used by the retained lines, namely the environments \texttt{lemma} on separate counters, together with the standard packages those lines need.  The retained environments print in this document as Lemma 1 in order of appearance, whereas the article prints the same items with the numbering of its own full document.  The complete native source of the article as distributed in the arXiv e-print for this exact version is bundled unchanged as \texttt{sources/208207/source0.tex}.
\begin{equation}\label{eq:K0-conditions}
K_0>e^4,\qquad
d_0:=\frac{1}{\sqrt{\log K_0}}\leq\eta^2,
\qquad
\eta\alpha_0d_0^2\leq1.
\end{equation}

\begin{lemma}
\label{lem:chord-lower-bound}
For every \(k\ge 0\),
\[
\|p^{k+1}-p^k\|
\ge \frac{\beta_k}{\pi}.
\]
\end{lemma}

\begin{lemma}
\label{lem:elementary}
Fix \(\eta\in(0,1/2]\) and choose \(K_0\) such that
\eqref{eq:K0-conditions} holds. Then there exist absolute constants
\(\widehat C_1,\widehat C_2,\widehat C_3,\widehat C_4>0\), independent
of \(k\), \(K_0\), and \(\eta\), such that
\begin{align}
    0<d_k-d_{k+1}
    &\leq \widehat C_1\alpha_kd_k^3,
    \label{eq:d-first}\\
    \norm{v_k}
    &\leq \widehat C_2\eta d_k^2,
    \label{eq:v-size}\\
    \norm{v_{k+1}-v_k}
    &\leq \widehat C_3\eta\,\norm{p^{k+1}-p^k},
    \label{eq:v-consecutive}\\
    0\leq a_k
    &\leq \widehat C_4\eta^2d_k^2.
    \label{eq:a-size}
\end{align}
\end{lemma}

\begin{proof}
Set
$
    u_k:=u(\varphi_k),
    \;
    e_k:=(-\sin\varphi_k,\cos\varphi_k),
    \;
    \beta_k:=\varphi_{k+1}-\varphi_k
             =\eta\alpha_kd_k^2.
$
Choose \(K_0\) such that
\eqref{eq:K0-conditions} holds. Then, for every \(k\geq0\),
\[
r_k\geq\frac12,
\qquad
0\leq\beta_k\leq1,
\qquad
d_k^2\leq\eta,
\qquad
d_k^3\leq\eta^2.
\label{eq:app-K-conditions}
\]
%Indeed, \(d_k\leq d_0\leq\eta^2\) and \(0<\eta\leq1\).


\noindent To prove \eqref{eq:d-first}, define
\[
d(t):=(\log t)^{-1/2},\qquad t\geq K_0.
\]
Since
\[
d'(t)
=-\frac{1}{2t(\log t)^{3/2}}
=-\frac{d(t)^3}{2t},
\]
and \(d_k=d(k+K_0)\), the mean-value theorem implies that, for some
\(\xi_k\in(k,k+1)\),
\begin{align*}
0<d_k-d_{k+1}
&=-d'(\xi_k+K_0)\\
&=\frac{1}{2(\xi_k+K_0)
[\log(\xi_k+K_0)]^{3/2}}\\
&\leq
\frac{1}{2(k+K_0)[\log(k+K_0)]^{3/2}}\\
&\leq
\frac{2}{2(k+K_0+1)[\log(k+K_0)]^{3/2}}\\
&=2\alpha_k d_k^3,
\end{align*}
where the second inequality follows from \(k+K_0+1\leq 2(k+K_0)\).
Thus, \eqref{eq:d-first} holds with
\(\widehat C_1=2\).


Note that
\begin{equation}
    \left|d_{k+1}^2-d_k^2\right|=\left|d_{k+1}-d_k\right|\left|d_{k+1}+d_k\right|\le 2d_k\left|d_{k+1}-d_k\right|
    \leq 4\alpha_kd_k^4.
    \label{eq:app-d-square}
\end{equation}

Define the scaled radial increment
\[
    \rho_k
    :=\frac{r_{k+1}-r_k}{\alpha_k}
    =\frac{d_k-d_{k+1}}{\alpha_k}.
\]
We have that
\begin{equation}
    0\leq \rho_k\leq 2d_k^3.
    \label{eq:app-rho-size}
\end{equation}
%Indeed,
%\begin{eqnarray*}
%\rho_k
%&=&2(k+K_0+1)
%\left(
%\frac{1}{\sqrt{\log(k+K_0)}}
%-\frac{1}{\sqrt{\log(k+K_0+1)}}
%\right)\\
%&\leq&
%2(k+K_0+1)
%\frac{\log\left(1+\frac{1}{k+K_0}\right)}
%     {[\log(k+K_0)]^{3/2}}\\
%&\leq&
%2(k+K_0+1)
%\frac{1}{(k+K_0)[\log(k+K_0)]^{3/2}}\\
%&\leq&
%\frac{4}{[\log(k+K_0)]^{3/2}}
%=4d_k^3,
%\end{eqnarray*}
%where the second-to-last inequality follows from $\log(1+x)\le x$ when $x\in[0,1)$.
Define
$
    F(t):=2(t+1)\bigl(d(t)-d(t+1)\bigr).$
Then \(\rho_k=F(k+K_0)\).  Since
\[
    d(t)-d(t+1)
    =\frac12\int_t^{t+1}
      \frac{ds}{s(\log s)^{3/2}},
\]
we can write
\[
    F(t)=(t+1)\int_t^{t+1}f(s)\,ds,
    \qquad
    f(s):=\frac{1}{s(\log s)^{3/2}}.
\]
Differentiation gives
\[
    F'(t)
    =\int_t^{t+1}\bigl(f(s)+(t+1)f'(s)\bigr)\,ds,
\]
where
\[
    f'(s)
    =-\frac{1}{s^2(\log s)^{3/2}}
     -\frac{3}{2s^2(\log s)^{5/2}}.
\]
Using the preceding representation of $F'(t)$, together with the mean-value theorem, there exists $\widehat s\in(t,t+1)$ such that
\begin{eqnarray*}
|F'(t)|
&\leq& \frac{1}{\widehat s(\log(\widehat s))^{3/2}}
+(t+1)\left(\frac{1}{\widehat s^2(\log(\widehat s))^{3/2}}
+\frac{3}{2\widehat s^2(\log(\widehat s))^{5/2}}\right)\nn\\
&\leq&\frac{1}{t(\log(t))^{3/2}}
+(t+1)\left(\frac{1}{t^2(\log(t))^{3/2}}
+\frac{3}{2t^2(\log(t))^{5/2}}\right)\nn\\
&\leq&\frac{3}{t(\log(t))^{3/2}}.
\end{eqnarray*}

Since \(t\ge e^{4}\), we have
$
\frac{t+1}{t}
=1+\frac1t
\le
\frac43,
\;
1+\frac{3}{2\log t}
\le
1+\frac38
=
\frac{11}{8}.$
Hence,
$
1+\frac{t+1}{t}
\left(
1+\frac{3}{2\log t}
\right)
\le
1+\frac43\cdot\frac{11}{8}
=
\frac{17}{6}
<
3,
$
which implies
$
|F'(t)|
\le
\frac{17}{6t(\log t)^{3/2}}
<
\frac{3}{t(\log t)^{3/2}}.
$
%where the last inequality holds for $t\ge e^4$. Indeed, since $\log(t)\ge 4$, we have
%\[
%\frac{1}{t^2(\log(t))^{5/2}}
%\leq
%\frac{1}{t^2(\log(t))^{3/2}}.
%\]
%Recalling that \(\rho_k=F(k+K_0)\) and applying the preceding estimate, we obtain
%\begin{equation}
%    |\rho_{k+1}-\rho_k|
%    \leq C'_3\alpha_kd_k^3.
%    \label{eq:app-rho-difference}
%\end{equation}

By the mean-value theorem, there exists
$\zeta_k\in(k+K_0,k+K_0+1)$ such that
\begin{eqnarray*}
 &&|\rho_{k+1}-\rho_k|
 =|F(k+K_0+1)-F(k+K_0)|
 =|F'(\zeta_k)|\nn\\
 &&
 \leq
 \frac{3}{\zeta_k(\log \zeta_k)^{3/2}}\leq\frac{3}{(k+K_0)(\log (k+K_0))^{3/2}}\leq\frac{6}{(k+K_0+1)(\log (k+K_0))^{3/2}}
 \leq 12\alpha_k d_k^3.
\end{eqnarray*}
%where we use
%\[
%\frac{1}{\zeta_k}\leq\frac{1}{k+K_0},
%\qquad
%\frac{1}{k+K_0}\leq\frac{2}{k+K_0+1}.
%\]
\noindent Next, since
\[
    u_{k+1}=\cos\beta_k\,u_k+\sin\beta_k\,e_k,
\]
we have
\begin{eqnarray}
    &&v_k=\frac{r_ku_k-r_{k+1}u_{k+1}}{\alpha_k}=A_ku_k-B_ke_k,\nn\\
    &&A_k:=-\rho_k+c_k,
    \;\;
    c_k:=\frac{r_{k+1}(1-\cos\beta_k)}{\alpha_k},
    \;\;
    B_k:=\frac{r_{k+1}\sin\beta_k}{\alpha_k}.
    \label{eq:app-v-frame}
\end{eqnarray}
%where
%\begin{equation}
%
%\end{equation}
Using \(1-\cos t\leq t^2/2\), \(|\sin t|\leq|t|\), and the definition of
\(\beta_k\), we obtain
\begin{equation}
    0\leq c_k\leq \frac{1}{2}\eta^2\alpha_kd_k^4,
    \qquad
    |B_k|\leq\eta d_k^2.
    \label{eq:app-c-B-size}
\end{equation}
Combining \eqref{eq:app-rho-size}, \eqref{eq:app-v-frame},
\eqref{eq:app-c-B-size}, and \(d_k\leq\eta/2\) and $\eta\alpha_k d_k^2\le 1$, we conclude that
\[
    \norm{v_k}
    \leq |A_k|+|B_k|
    \leq 2 d_k^3+\frac{1}{2}\eta^2 \alpha_k d_k^4 +\eta d_k^2
    \leq 2 \cdot\frac{\eta}{2}\cdot d_k^2+\frac{1}{2}\eta d_k^2+\eta d_k^2\le\frac{5}{2}\eta d_k^2.
\]
This proves \eqref{eq:v-size} by setting $\widehat C_2:=5/2$.

Next, we show (\ref{eq:v-consecutive}).
First, monotonicity of \(\alpha_k\) and \(d_k\) gives
\begin{equation}
    |c_{k+1}-c_k|
    \leq c_{k+1}+c_k
    \leq \eta^2\alpha_kd_k^4.
    \label{eq:app-c-difference}
\end{equation}
 Moreover,
\eqref{eq:d-first} and \eqref{eq:app-d-square} imply
\[
    |r_{k+2}-r_{k+1}|\leq 2 \alpha_{k+1}d_{k+1}^3
\]
and
\begin{eqnarray}\label{addkey}
    \left|
       r_{k+2}d_{k+1}^2-r_{k+1}d_k^2
    \right|
    \leq 6\alpha_kd_k^4.
\end{eqnarray}
%Consequently,
%\begin{align}
%    |B_{k+1}-B_k|
%    &\leq
%    \eta\left|
%       r_{k+2}d_{k+1}^2-r_{k+1}d_k^2
%    \right|
%    +C_7\eta d_k^2(\beta_k^2+\beta_{k+1}^2)\leq C_8\eta\alpha_kd_k^4.
%    \label{eq:app-B-difference}
%\end{align}
%Here all higher-order terms are absorbed using
%\(0<\alpha_k,d_k,\eta\leq1\).


We next derive an upper bound on $|B_{k+1}-B_k|$. Recall that
\[
B_k=\eta r_{k+1}d_k^2\sinc(\beta_k),
\qquad
B_{k+1}=\eta r_{k+2}d_{k+1}^2\sinc(\beta_{k+1}),
\]
where
\[
\sinc(t):=
\begin{cases}
\dfrac{\sin t}{t}, & t\neq 0,\\[1mm]
1, & t=0.
\end{cases}
\]
Hence
\begin{align*}
B_{k+1}-B_k
={}&
\eta r_{k+2}d_{k+1}^2\sinc(\beta_{k+1})
-\eta r_{k+1}d_k^2\sinc(\beta_k)\\
={}&
\eta\bigl(r_{k+2}d_{k+1}^2-r_{k+1}d_k^2\bigr)
+\eta r_{k+2}d_{k+1}^2
   \bigl(\sinc(\beta_{k+1})-1\bigr)
-\eta r_{k+1}d_k^2
   \bigl(\sinc(\beta_k)-1\bigr).
\end{align*}
Therefore,
\begin{align*}
|B_{k+1}-B_k|
\le{}&
\eta\left|r_{k+2}d_{k+1}^2-r_{k+1}d_k^2\right|+
\eta r_{k+2}d_{k+1}^2
\left|\sinc(\beta_{k+1})-1\right|+
\eta r_{k+1}d_k^2
\left|\sinc(\beta_k)-1\right|.
\end{align*}
Since \(0\le r_k\le 1\), \(d_{k+1}\le d_k\), and, for \(|t|\le 1\),
$
|\sinc(t)-1|\le \frac{t^2}{6},
$
we obtain
\begin{align*}
&\eta r_{k+2}d_{k+1}^2
\left|\sinc(\beta_{k+1})-1\right|
+
\eta r_{k+1}d_k^2
\left|\sinc(\beta_k)-1\right|\\
&\qquad\le
\frac{\eta}{6}
\left(
d_{k+1}^2\beta_{k+1}^2+d_k^2\beta_k^2
\right)\\
&\qquad\le
\frac{\eta d_k^2}{6}
\left(
\beta_{k+1}^2+\beta_k^2
\right).
\end{align*}
Thus,
\[
|B_{k+1}-B_k|
\le
\eta\left|r_{k+2}d_{k+1}^2-r_{k+1}d_k^2\right|
+
\frac{1}{6}\eta d_k^2
\left(\beta_k^2+\beta_{k+1}^2\right).
\]

Using (\ref{addkey})
%\[
%\left|r_{k+2}d_{k+1}^2-r_{k+1}d_k^2\right|
%\le 6\alpha_kd_k^4,
%\]
together with
\[
\beta_k=\eta\alpha_kd_k^2,
\qquad
\beta_{k+1}\le\beta_k,
\]
we further have
\begin{align*}
\frac{1}{6}\eta d_k^2(\beta_k^2+\beta_{k+1}^2)
&\le
\frac{1}{3}\eta d_k^2\beta_k^2=
\frac{1}{3}\eta^3\alpha_k^2d_k^6=
\frac{1}{3}(\eta^2\alpha_kd_k^2)\,
\eta\alpha_kd_k^4\le
\frac{1}{3}\eta\alpha_kd_k^4,
\end{align*}
where the last inequality uses
\(0<\eta\alpha_kd_k^2\le1\). Consequently,
\[
|B_{k+1}-B_k|
\le
\frac{19}{3}\eta\alpha_kd_k^4.
\]



Equations \eqref{eq:app-rho-size}, \eqref{eq:app-c-B-size} and
\eqref{eq:app-c-difference}, together with the preceding estimate for $|\rho_{k+1}-\rho_k|$, yield
\begin{align}
    |A_k|
    &\leq 2 d_k^3+\frac{\eta^2}{2}\alpha_kd_k^4,
    \label{eq:app-A-size}\\
    |A_{k+1}-A_k|
    &\leq 12\alpha_kd_k^3
          +\eta^2\alpha_kd_k^4.
    \label{eq:app-A-difference}
\end{align}
The orthonormal frames \((u_k,e_k)\) and
\((u_{k+1},e_{k+1})\) differ by a rotation through \(\beta_k\); hence
\[
    \norm{u_{k+1}-u_k}=2\sin(\beta_k/2)\leq\beta_k,
    \qquad
    \norm{e_{k+1}-e_k}=2\sin(\beta_k/2)\leq\beta_k.
\]
Using the representation \eqref{eq:app-v-frame}, we therefore obtain
\begin{align}
    \norm{v_{k+1}-v_k}
    &\leq |A_{k+1}-A_k|+|B_{k+1}-B_k|\notag\\
    &+\bigl(|A_{k+1}|+|B_{k+1}|\bigr)\beta_k\notag\\
    &\leq 12 \alpha_k d_k^3+\left(2\eta^2\alpha_k+\frac{19}{3}\eta\alpha_k\right)d_k^4+2\eta\alpha_k d_k^5+\frac{1}{2}\eta^3\alpha_k^2d_k^6.
    \label{eq:app-v-difference-preliminary}
\end{align}

Indeed, dividing the right-hand side of \eqref{eq:app-v-difference-preliminary} by
$\eta^2\alpha_kd_k^2$ and using
$d_k\leq\eta^2$, $\eta\leq 1/2$, $\alpha_k\leq 1$, and
$\eta\alpha_kd_k^2\leq 1$, we obtain
\[
    12+2\eta^4+\frac{19}{3}\eta^3+2\eta^5
    +\frac12\eta^9
    <\frac{27}{2},
\]
where the last inequality follows from $\eta\leq 1/2$.
Thus
\begin{equation}
    \norm{v_{k+1}-v_k}
    \leq \frac{27}{2}\eta^2\alpha_kd_k^2.
    \label{eq:app-v-difference-final}
\end{equation}
On the other hand, Lemma~\ref{lem:chord-lower-bound} gives
\begin{equation}
    \norm{p^{k+1}-p^k}
    \geq\frac{\beta_k}{\pi}
    =\frac{\eta\alpha_kd_k^2}{\pi}.
    \label{eq:app-step-lower-bound}
\end{equation}
Combining \eqref{eq:app-v-difference-final} and
\eqref{eq:app-step-lower-bound} proves
\[
    \norm{v_{k+1}-v_k}
    \leq \frac{27}{2}\pi\eta\norm{p^{k+1}-p^k},
\]
which is \eqref{eq:v-consecutive} by setting $\widehat C_3:= \frac{27}{2}\pi$.

\medskip
Next, we prove \eqref{eq:a-size}.
Using \eqref{eq:v-size},
\begin{align*}
    a_k=\sum_{\ell=k}^{\infty}\alpha_\ell\norm{v_\ell}^2\leq (5/2)^2\eta^2
      \sum_{\ell=k}^{\infty}
      \frac{1}{2(\ell+K_0+1)\log^2(\ell+K_0)}.
\end{align*}
The integral test gives
\[
    \sum_{\ell=k}^{\infty}
      \frac{1}{(\ell+K_0+1)\log^2(\ell+K_0)}
      \leq\sum_{n=k+K_0}^{\infty}\frac{1}{n\log^2 n}
    \leq\frac{2}{\log(k+K_0)}
    =2 d_k^2.
\]
Therefore
\[
    0\leq a_k\leq \frac{25}{4} \eta^2d_k^2,
\]
which is \eqref{eq:a-size} by setting $\widehat C_4:=\frac{25}{4} $.
\end{proof}

\end{document}
